\clearpage\fancyhead[L]{\small KOMET.}\hypertarget{native-390-lower}{}\label{native-390-lower}
\begin{center}\textbf{KOMET.}\par\vspace{4pt}\hrule\vspace{4pt}\fontsize{10}{12}\selectfont Göttingische gelehrte Anzeigen. Stück 105. Seite 1041..1043. 3. Juli 1815.\par\vspace{4pt}\hrule\vspace{4pt}\end{center}

\par\indent Im 55. Stück dieser Zeitschrift [S. 382 d. B.] wurden die ersten Beobachtungen des jetzt sichtbaren Kometen von der hiesigen Sternwarte mitgeteilt. Wir können nun die folgenden zwei von Herrn Prof. \textsc{Gauss} hinzufügen, den zufällige Umstände zu einer recht langen Unterbrechung gezwungen hatten.

\begin{center}\fontsize{10}{12}\selectfont\setlength{\tabcolsep}{4pt}\renewcommand{\arraystretch}{1.15}\begin{tabular}{l|c|c|c}
1815 & Mittlere Zeit & Rektaszension & Deklination.\\\hline
April 2 & $9^{\mathrm u}\;11^{\mathrm m}\;55^{\mathrm s}$ & $61^\circ\;2'\;27''$ & $45^\circ\;39'\;57''\;\mathrm{N.}$\\
Juni 12 & $10\;53\;31$ & $165\;3\;3$ & $52\;8\;25$\\
\end{tabular}\end{center}

\Needspace{23\baselineskip}\par\indent Von den zahlreichen Beobachtungen, die auswärtige Astronomen mitgeteilt haben, führen wir hier nur die an, die Herr \textsc{Santini} an der Sternwarte in Padua angestellt hat. Der Sekretär der dortigen Sektion des K. K. Instituts, Herr Ritter \textsc{Brera}, hat sie Herrn Prof. \textsc{Gauss} in einem Schreiben vom 19. Mai zugesandt.

\SourceLeaf{391}\begin{center}\fontsize{10}{12}\selectfont\setlength{\tabcolsep}{5pt}\renewcommand{\arraystretch}{1.15}\begin{tabular}{l|c|c|c}
1815 & \shortstack{Mittlere Zeit\\in Padua} & \shortstack{Rektaszension} & \shortstack{Nördliche\\Deklination}\\\hline
April 24 & $9^{\mathrm u}\;39^{\mathrm m}\;50^{\mathrm s}$ & $80^\circ\;49'\;33''$ & $56^\circ\;3'\;16''$\\
28 & $8\;57\;49$ & $85\;57\;54$ & $57\;37\;53$\\
Mai 1 & $8\;57\;22$ & $90\;16\;13$ & $58\;41\;7$\\
2 & $9\;8\;51$ & $91\;47\;15$ & $59\;0\;32$\\
2 & $9\;40\;22$ & $91\;48\;48$ & $59\;0\;40$\\
6 & $9\;29\;48$ & $98\;13\;9:$ & $\text{wolkig}$\\
8 & $9\;10\;55$ & $101\;40\;4$ & $60\;35\;36$\\
8 & $9\;45\;26$ & $101\;42\;13$ & $60\;37\;0$\\
10 & $9\;9\;11$ & $105\;14\;41$ & $60\;59\;45:$\\
11 & $8\;37\;5$ & $107\;4\;0$ & $61\;6\;4$\\
11 & $9\;22\;30$ & $107\;7\;18$ & $61\;7\;30$\\
12 & $9\;3\;20$ & $108\;57\;37$ & $61\;16\;32:$\\
12 & $9\;48\;15$ & $109\;2\;11$ & $61\;16\;7$\\
17 & $10\;15\;15$ & $118\;47\;45$ & $61\;30\;45$\\
17 & $10\;31\;5$ & $118\;49\;35$ & $61\;30\;47$\\
\end{tabular}\end{center}

\par\indent Diese Beobachtungen wurden mit einem Kreismikrometer am \textsc{Adams}schen Quadranten angestellt.

\par\indent Mehrere Astronomen haben inzwischen neue parabolische Elemente bestimmt, die alle nur wenig von den ersten Elementen abweichen, die Herr Prof. \textsc{Gauss} bestimmt hat und die an der genannten Stelle mitgeteilt wurden. So gut alle diese Elemente die früheren Beobachtungen auch wiedergaben, sie teilten doch das gleiche Schicksal. Die späteren Beobachtungen wichen schneller und stärker davon ab, als man es sonst bei solchen Berechnungen gewohnt ist. Das deutete bereits auf eine merkliche Abweichung der Bahn von einer Parabel hin. Da der Komet nun immer lichtschwächer und seine Beobachtung immer schwieriger wird, hielt Herr Prof. \textsc{Gauss} es für nützlich, die parabolischen Elemente noch einmal zu korrigieren, um die weiteren Beobachtungen zu erleichtern. Aus den Beobachtungen vom 6. März und 25. April (beide von Herrn Dr. \textsc{Olbers}) und vom 12. Juni erhielt er folgende Ergebnisse:

\Needspace{12\baselineskip}\begin{center}\fontsize{10}{12}\selectfont Durchgang durch das Perihel am 25. April.\par $11^{\mathrm u}\;41^{\mathrm m}\;19^{\mathrm s}$ M. Z. in Göttingen\end{center}
\begin{center}\fontsize{10}{12}\selectfont\setlength{\tabcolsep}{3pt}\renewcommand{\arraystretch}{1.15}\begin{tabular}{lcr}
Länge des Perihels & $.\;.\;.\;.\;.\;.\;.\;.\;.\;.\;.\;.$ & $147^\circ\;35'\;55''$\\
Kleinster Abstand & $.\;.\;.\;.\;.\;.\;.\;.\;.\;.\;.\;.\;.\;.$ & $1.23024$\\
Aufsteigender Knoten & $.\;.\;.\;.\;.\;.\;.\;.\;.\;.\;.\;.\;.$ & $82^\circ\;43'\;6''$\\
Neigung der Bahn & $.\;.\;.\;.\;.\;.\;.\;.\;.\;.\;.\;.\;.$ & $44\;43\;13$\\
\end{tabular}\end{center}\begin{center}Bewegung rechtläufig.\end{center}

\par\indent Diese Elemente kommen den drei Beobachtungen in jeder Hinsicht sehr nahe. Dennoch weichen sie sehr stark von den dazwischenliegenden Beobachtungen ab (bei der Rektaszension um \emph{bis zu elf Minuten}). Diese Erscheinung kann oft vorkommen und sollte nicht überraschen. Daraus folgt einerseits, dass die Bahn beträchtlich von einer Parabel abweicht, und andererseits, dass sich die drei genannten Beobachtungen zur Bestimmung des Kegelschnitts nicht eignen, da sie ihn für sich genommen bis zu einem gewissen Grad unbestimmt lassen. (Der eigentliche Grund für dieses Phänomen liegt darin, dass die Richtung der geozentrischen Bewegung um die Zeit der mittleren Beobachtung fast mit der Richtung eines vom Kometen zur Sonne gezogenen Großkreises zusammenfiel. Siehe die letzten Artikel im ersten Abschnitt des zweiten Buchs der \emph{Theor. Mot. C. C.}.) Es ist daher zu erwarten, dass auch die folgende kurze, aus diesen parabolischen Elementen berechnete Ephemeride von Anfang an merklich von der tatsächlichen Bewegung abweicht. Sie wird zu kleine Rektaszensionen und zu große Deklinationen liefern, dürfte aber doch genau genug sein, um den Kometen leichter zu finden:

\SourceLeaf{392}\begin{center}\fontsize{10}{12}\selectfont\setlength{\tabcolsep}{5pt}\renewcommand{\arraystretch}{1.15}\begin{tabular}{l|c|c}
\shortstack{$10^{\mathrm u}\;53^{\mathrm m}$\\in Göttingen} & \shortstack{Rektaszension} & Deklination\\\hline
Juni 20 & $174^\circ\;46'$ & $46^\circ\;47'\;\mathrm{N.}$\\
24 & $178\;55$ & $43\;56$\\
28 & $182\;41$ & $41\;3$\\
Juli 2 & $186\;7$ & $38\;8$\\
6 & $189\;16$ & $35\;16$\\
10 & $192\;11$ & $32\;25$\\
14 & $194\;53$ & $29\;40$\\
\end{tabular}\end{center}

\par\indent Am 20. Juni beträgt die Lichtstärke noch $\frac{3}{10}$ der Lichtstärke am Tag der Entdeckung des Kometen, am 14. Juli hingegen nur die Hälfte davon.

\par\indent Obwohl es für eine genaue Bestimmung der wahren Bahn des Kometen noch zu früh ist, hielt Herr Prof. \textsc{Gauss} eine vorläufige Bestimmung für interessant. Sie gibt wenigstens eine ungefähre Vorstellung davon, wie unerwartet stark die Bahn von der Parabel abweicht, und kann außerdem dazu dienen, die Beobachtungen im Hinblick auf eine genauere Bestimmung zu erörtern und vorzubereiten. Dem Zweck entsprechend wurde die Rechnung nur überschlägig ausgeführt, allein auf fünf vereinzelte Beobachtungen gestützt und ohne jede Berücksichtigung von Präzession, Aberration, Nutation und Parallaxe. Diese Umstände peinlich genau zu berücksichtigen, wäre bei einer solchen ersten Bestimmung verlorene Mühe gewesen.

\Needspace{13\baselineskip}\begin{center}
Vorläufige \emph{elliptische} Elemente des Kometen von 1815.
\end{center}

\begin{center}\fontsize{10}{12}\selectfont Durchgang durch das Perihel am 26. April $0^{\mathrm u}\;21^{\mathrm m}\;25^{\mathrm s}$\end{center}
\begin{center}\fontsize{10}{12}\selectfont\setlength{\tabcolsep}{3pt}\renewcommand{\arraystretch}{1.15}\begin{tabular}{lcr}
Länge des Perihels & $.\;.\;.\;.\;.\;.\;.\;.\;.\;.\;.$ & $148^\circ\;58'\;48''$\\
Abstand im Perihel & $.\;.\;.\;.\;.\;.\;.\;.\;.$ & $1.21349$\\
Aufsteigender Knoten & $.\;.\;.\;.\;.\;.\;.\;.\;.\;.\;.$ & $83\;26\;21$\\
Neigung der Bahn & $.\;.\;.\;.\;.\;.\;.\;.\;.\;.\;.\;.\;.$ & $44\;30\;43$\\
Exzentrizität & $.\;.\;.\;.\;.\;.\;.\;.\;.\;.\;.\;.\;.\;.$ & $0.933149$\\
\end{tabular}\end{center}

\par\indent Daraus ergeben sich ferner ein Abstand im Aphel von 35.091 und eine Umlaufzeit von $77\frac{1}{8}$ Jahren, fast so lang wie die des \textsc{Halley}schen Kometen. Mit diesen Elementen entfallen alle jene groben Abweichungen der parabolischen Bahn vollständig. Herr Prof. \textsc{Gauss} glaubt sie daher schon jetzt als eine Annäherung an die Wahrheit ansehen zu dürfen. Selbst wenn eine spätere genauere Rechnung eine merklich längere Umlaufzeit ergibt, wird sie doch kaum mehr als 100 Jahre betragen können. Man kann diesen Kometen daher schon jetzt als eine der bemerkenswertesten Erscheinungen in den Annalen der Astronomie ansehen.
